\documentclass[11pt]{amsart}
\usepackage{amssymb}
\usepackage[margin=1.2in]{geometry}
\usepackage{longtable}

\newtheorem{theorem}{Theorem}[section]
\newtheorem{lemma}[theorem]{Lemma}
\newtheorem{proposition}[theorem]{Proposition}
\newtheorem{corollary}[theorem]{Corollary}
\theoremstyle{definition}
\newtheorem{definition}[theorem]{Definition}
\theoremstyle{remark}
\newtheorem{remark}[theorem]{Remark}

\newcommand{\R}{R}
\newcommand{\Pn}{\mathbb{P}^n}
\newcommand{\M}{M_N(R)}
\newcommand{\AZ}{\mathrm{AZ}_n}
\newcommand{\SB}{\mathrm{SB}_n}
\newcommand{\RI}{\mathrm{RI}}
\newcommand{\PGL}{\mathrm{PGL}_N(R)}
\newcommand{\Aut}{\mathrm{Aut}}
\newcommand{\tr}[1]{\Vert #1\Vert}
\newcommand{\trT}[1]{\Vert #1\Vert_T}
\newcommand{\Col}{\mathrm{Col}}
\newcommand{\Roots}{\mathrm{Roots}}

\title{Ch\^atelet's theorem in synthetic algebraic geometry:\\ a checked proof}
\author{Thierry Coquand}
\date{October 23, 2026}

% stamp on every page: file name and date/time of compilation
\newcount\stamphour \newcount\stampmin
\stamphour=\time \divide\stamphour by 60
\stampmin=\time \multiply\stamphour by 60 \advance\stampmin by -\stamphour \divide\stamphour by 60
\edef\stamptime{\the\year-\ifnum\month<10 0\fi\the\month-\ifnum\day<10 0\fi\the\day\space
 \ifnum\stamphour<10 0\fi\the\stamphour:\ifnum\stampmin<10 0\fi\the\stampmin}
\AddToHook{shipout/foreground}{\put(0.5\paperwidth,-\dimexpr\paperheight-0.45in\relax)%
 {\makebox[0pt]{\scriptsize\ttfamily\jobname.pdf\quad\stamptime}}}
\begin{document}

\begin{abstract}
We present the proof of Ch\^atelet's theorem in synthetic algebraic
geometry that has been checked in the proof assistant Unit: a
Severi--Brauer variety with a point is projective space. The proof follows
the note \emph{Ch\^atelet's theorem in synthetic algebraic geometry}
(sevb), with three changes. For the \'etale topology it is checked from
the axioms of synthetic algebraic geometry and a left exact
sheafification; the descent of finite freeness that it needs is proved in
a companion note. Right ideals are required to be free
\emph{direct summands}. All ``same dimension, hence equal'' steps are
replaced by one fact on square matrices over a local ring, proved without
determinants. Compactness of the spectrum of $R[X]/(f)$, for $f$ monic, is
proved algebraically. We also present the part of sevb's Section~1 that has
been checked: for the topology generated by roots of monic polynomials,
$R$ is a sheaf and so is every scheme. The last section lists, for each
statement, the corresponding checked definition.
\end{abstract}

\maketitle

\section{Setting and statement}

$R$ is the ring of synthetic algebraic geometry: it is local (that is, $1\neq
0$, and if $a+b$ is invertible then $a$ or $b$ is invertible), and it
satisfies synthetic quasi-coherence and Zariski-local choice
\cite{CCH24}. We work in univalent type theory; ``merely'' means inside the
propositional truncation $\tr{-}$. Fix $n\geq 1$ and put $N=n+1$.

\emph{Projective space} $\Pn$ is the type of vectors $x\in R^N$ with some
coordinate invertible, modulo $x\sim ux$ for $u$ a unit. $\M$ is the algebra
of $N\times N$ matrices, $E_{ij}$ the elementary matrices, $\PGL$ the
invertible linear maps of $R^N$ modulo unit scalars.

\subsection*{The topology}
A topology $T$ is given by a family of propositions $G_i$ ($i\in I$). A type
$X$ is a \emph{sheaf} if for every $i$ the constant map $X\to X^{G_i}$,
$x\mapsto(\_\mapsto x)$, is an equivalence. We assume:
\begin{itemize}
\item[(H1)] $R$ is a sheaf.
\item[(H2)] For every type $A$ there is a proposition $\trT{A}$ that is a
  sheaf, with a map $A\to\trT A$, such that every map from $A$ to a
  proposition that is a sheaf factors through $\trT A$.
\item[(H3)] (Lexness) The type of sheaves $\Sigma_{X} \,(X\text{ is a sheaf})$
  is a sheaf.
\item[(H4)] (Descent) For an algebra $A$ whose carrier is a sheaf and a
  right ideal $I$ of $A$ (Definition~\ref{def:ri}), if $I$ has a basis
  $T$-locally then it merely has a basis.
\end{itemize}
For the \'etale topology of Section~\ref{sec:etale}, (H1) is proved there,
and (H4) is proved in \cite{descent} for the algebras that are used:
those with $\trT{A=\M}$, including $\M$ itself.

\begin{definition}
$\AZ$ is the type of $R$-algebras $A$ whose carrier is a sheaf and such that
$\trT{A=\M}$. $\SB$ is the type of sheaves $X$ such that $\trT{X=\Pn}$.
\end{definition}

\begin{theorem}[Ch\^atelet]\label{thm:main}
For $X:\SB$, we have $\tr X\to\tr{X=\Pn}$.
\end{theorem}

The proof has three parts. The map $A\mapsto\RI(A)$, right ideals of
rank~$N$, is an equivalence $\AZ\simeq\SB$
(Proposition~\ref{prop:331}); both types are deloopings of $\PGL$, since
$\Aut(\M)=\PGL=\Aut(\Pn)$. A point of $\RI(A)$ is a right ideal $I$, and
$A$ acts on $I$ by an isomorphism $A\simeq\M$ (Lemma~\ref{lem:341}). Hence
$X=\RI(A)=\RI(\M)=\Pn$ (Lemma~\ref{lem:delta}).

\section{Sheaves}

\begin{lemma}\label{lem:closure}
Sheaves are closed under $\Pi$ (over any type), $\Sigma$, identity types and
retracts. A proposition $P$ is a sheaf if and only if every map $G_i\to P$
has a value, i.e.\ $(G_i\to P)\to P$ for all $i$.
\end{lemma}
\begin{proof}
For $\Pi$ the inverse is pointwise. For $\Sigma_{a:A}B(a)$, the constant map
is the composite
\[\textstyle\Sigma_{a}B(a)\to\Sigma_{a}B(a)^{G}\to\Sigma_{\alpha:G\to
A}\Pi_{g}B(\alpha g)\cong(G\to\Sigma_aB(a)),\]
of a fibrewise equivalence, the base change of a family along the
equivalence $A\to A^G$ (an equivalence by induction on equivalences), and a
canonical isomorphism. For $x,y:X$ the constant map $(x=y)\to(G\to x=y)$ is
$\mathrm{ap}$ of the equivalence $X\to X^G$ followed by function
extensionality. For a retract $s:B\to A$, $r:A\to B$, $rs=1$, with $A$ a sheaf, the
inverse of the constant map of $B$ sends $h:G\to B$ to $r(a)$, where $a$
is the value of $s\circ h$. For a
proposition both composites of the two maps are equalities between
elements of propositions.
\end{proof}

\begin{lemma}\label{lem:propT}
The type $\mathrm{Prop}_T$ of propositions $P$ with sections $(G_i\to P)\to
P$ is a sheaf, and two elements of it with the same inhabitants are equal.
\end{lemma}
\begin{proof}
By Lemma~\ref{lem:closure} it is a retract of $\Sigma_{X\text{
sheaf}}(X\text{ is a proposition})$. This is a $\Sigma$ over the type of
sheaves, a sheaf by (H3), of a family of sheaves, since ``$X$ is a
proposition'' is a $\Pi$ of identity types. The second claim is
univalence; the sections form a proposition.
\end{proof}

\begin{lemma}[Delooping, sevb 3.1.2]\label{lem:deloop}
Let $A,B$ be sheaves with base points $a_0,b_0$, each \emph{connected}
($\trT{x=y}$ for all $x,y$), and $f:A\to B$ with $p:f(a_0)=b_0$. If
$\Omega f:(a_0=a_0)\to(b_0=b_0)$, $q\mapsto p^{-1}\cdot f(q)\cdot p$, is an
equivalence, then $f$ is an equivalence.
\end{lemma}
\begin{proof}
``$\mathrm{ap}_f:(x=y)\to(fx=fy)$ has contractible fibres'' is a proposition
and a sheaf (Lemma~\ref{lem:closure}). By connectedness we may take
$x=y=a_0$, where $\mathrm{ap}_f$ is $\Omega f$ conjugated by $p$. So the
fibres of $f$ are propositions; they are sheaves, so by connectedness of
$B$ each is inhabited, as the fibre over $b_0$ is.
\end{proof}

\section{Square matrices over a local ring}

\begin{lemma}[K1]\label{lem:K1}
If $\sum_i\lambda_iv_i=1$ in $R$, then merely some $v_i$ is invertible.
\end{lemma}
\begin{proof}
By locality some $\lambda_iv_i$ is invertible, hence $v_i$ is.
\end{proof}

\begin{lemma}[K2]\label{lem:K2}
If $L,M\in M_m(R)$ and $LM=1$ then $ML=1$.
\end{lemma}
\begin{proof}
Induction on $m$; the case $m=0$ is trivial. Let $t$ be the last index and
$c$ the last column of $M$. Since $\sum_kL_{tk}c_k=1$, by locality $c_t$ is
invertible or merely some $c_k$, $k\neq t$, is. In the second case
$E=1+aE_{tk}$, with $a=(1-c_t)c_k^{-1}$, satisfies $(1-aE_{tk})E=1$ and
$(EM)_{tt}=1$. So merely there is $P$ with a left inverse and $(PM)_{tt}$
invertible. Scaling row $t$ by $(PM)_{tt}^{-1}$ and clearing the rest of
the last column is a matrix $F$ with an explicit left inverse. So $Q=FP$
has a left inverse $Q'$ and the last column of $QM$ is $e_t$. The claim is
a proposition, and it is invariant under $(L,M)\mapsto(LQ',QM)$: if
$(QM)(LQ')=1$ then $ML=Q'Q\,ML\,Q'Q=Q'(QM\,LQ')Q=1$. So we may assume
$M=\begin{pmatrix}A&0\\ r&1\end{pmatrix}$, and write
$L=\begin{pmatrix}B&s\\ u&d\end{pmatrix}$. Then $LM=1$ says $s=0$, $d=1$,
$BA=1$ and $uA=-r$. By induction $AB=1$, so $rB=-uAB=-u$, and
$ML=\begin{pmatrix}AB&As\\ rB+u&rs+d\end{pmatrix}=1$.
\end{proof}

\section{Skolem--Noether}

\begin{proposition}[Skolem--Noether]\label{prop:sn}
Every unital algebra endomorphism $\varphi$ of $\M$ is merely of the form
$X\mapsto WXW^{-1}$.
\end{proposition}
\begin{proof}
Put $e_{ab}=\varphi(E_{ab})$, so $e_{ab}e_{cd}=\delta_{bc}e_{ad}$ and
$\sum_je_{jj}=1$. Fix an index $t$. Since $\sum_j(e_{jj})_{tt}=1$, by
Lemma~\ref{lem:K1} merely $c=(e_{jj})_{tt}$ is invertible for some $j$. Put
$W_{xy}=(e_{yj})_{xt}$ and $L_{xy}=c^{-1}(e_{jx})_{ty}$. Then
$(LW)_{xy}=c^{-1}(e_{jx}e_{yj})_{tt}=\delta_{xy}$, so $WL=1$ by
Lemma~\ref{lem:K2}. Both $\varphi(X)W$ and $WX$ have entries
$\varphi(XE_{yj})_{xt}$, since $XE_{yj}=\sum_zX_{zy}E_{zj}$. Hence
$\varphi(X)=WXL$.
\end{proof}

\begin{corollary}\label{cor:alpha}
The map $\alpha:\PGL\to\Aut(\M)$, $[G]\mapsto(X\mapsto GXG^{-1})$, is an
equivalence.
\end{corollary}
\begin{proof}
It is well defined since scalars cancel. It is surjective by
Proposition~\ref{prop:sn}. If $GXG^{-1}=G'XG'^{-1}$ for all $X$, then
$P=G'^{-1}G$ commutes with every $E_{ij}$, so $P=\mu\cdot1$ with $\mu$ a
unit, and $[G]=[G']$.
\end{proof}

We also use the equivalence $\beta:\PGL\to\Aut(\Pn)$, the action of
$\PGL$ on $\Pn$: this is the description of the automorphisms of projective
space of \cite{Che+24} (Corollary 5.6 there), checked earlier for every
$n\geq1$.

\section{Right ideals and the map $\delta$}

\begin{definition}\label{def:ri}
Let $A$ be an $R$-algebra. A \emph{right ideal} of $A$ is a map
$I:A\to\mathrm{Prop}_T$ closed under $0$, $+$, scalars and right
multiplication by $A$. A \emph{basis} of $I$ is a linear bijection $f:R^N\to
I$ together with a linear map $g:A\to R^N$, defined on all of $A$, with
$gf=1$. $\RI(A)$ is the type of right ideals $I$ with $\trT{I\text{ has a
basis}}$.
\end{definition}

\begin{remark}
The left inverse $g$ makes $I$ a free direct summand of $A$, as in sevb's
Grassmannian. It is needed: for $t$ a non-invertible non-zero-divisor,
$(t)\subseteq R$ is a right ideal of $M_1(R)=R$ that is free of rank one,
but it is not of the form $\delta(q)$ below.
\end{remark}

$\RI(A)$ is a sheaf by Lemmas~\ref{lem:closure} and~\ref{lem:propT}, for any
$A$.

For $x\in R^N$ and $X\in\M$ write $\Col_x(X)$ for: $X_{ij}x_k=X_{kj}x_i$ for
all $i,j,k$, i.e.\ every column of $X$ is proportional to $x$. If $x_k$ is
invertible, $\Col_x(X)$ holds if and only if $X=xy^T$ for some $y$, and
then $y=x_k^{-1}X_{k-}$.

\begin{lemma}[sevb 2.4.2]\label{lem:delta}
The map $\delta:\Pn\to\RI(\M)$, given by
\[\delta(q)=\{X\mid\Col_x(X)\text{ for every representative }x\text{ of }q\},\]
is an equivalence. On a representative, $\delta[x]=\{xy^T\mid y\in R^N\}$.
\end{lemma}
\begin{proof}
$\Col_x(X)$ is a sheaf proposition by (H1), and it is invariant under
$x\mapsto ux$; so membership in $\delta(q)$ needs no choice of
representative. $\Col_x$ is closed under the ideal operations, e.g.\
$(XY)_{ij}x_k=\sum_mX_{im}x_kY_{mj}=\sum_mX_{km}x_iY_{mj}=(XY)_{kj}x_i$. For
$x_k$ invertible, $f(y)=xy^T$ and $g(X)=x_k^{-1}X_{k-}$ form a basis of
$\delta[x]$.

\emph{Injective.} If $\delta(q)\subseteq\delta(q')$, with $q=[x]$,
$q'=[x']$, then $x\mathbf 1^T\in\delta[x']$ gives $x_ix'_k=x_kx'_i$, so $x,x'$
are proportional and $q=q'$.

\emph{Surjective.} Let $I:\RI(\M)$; the fibre of $\delta$ over $I$ is a
proposition, so by (H4) we may take a basis $(f,g)$ of $I$. Let $t$ be an
index and $M=f(e_t)$. Since $g_t$ is linear,
$1=g_t(M)=\sum_{a,b}M_{ab}\,g_t(E_{ab})$, so by Lemma~\ref{lem:K1} merely
some $M_{ab}$ is invertible. Let $c$ be the column $b$ of $M$. Then
$cy^T=M(e_by^T)\in I$ for all $y$. The matrices $M'_{ij}=g(ce_j^T)_i$ and
$L_{ji}=c_a^{-1}f(e_i)_{aj}$ satisfy $LM'=1$, since
$\sum_if(e_i)\,g(ce_{j'}^T)_i=f(g(ce_{j'}^T))=ce_{j'}^T$. By
Lemma~\ref{lem:K2}, $M'L=1$, which says $g(cr_i^T)=e_i$ with $(r_i)_j=L_{ji}$.
Since $cr_i^T\in I$ and $fg$ is the identity on $I$, $f(e_i)=cr_i^T$. So every element of $I$ is of the form
$cy^T$, and $I=\delta[c]$.
\end{proof}

In particular $\Pn$ is a sheaf. If $A:\AZ$ then $\RI(A):\SB$
(sevb~2.4.3): ``$\RI(A)=\Pn$'' follows from $A=\M$ and the lemma, so
$\trT{A=\M}$ gives $\trT{\RI(A)=\Pn}$.

\section{The equivalence $\AZ\simeq\SB$}

An algebra isomorphism $\varphi:A\simeq B$ acts on right ideals by
$\varphi[I]=\{y\mid\varphi^{-1}(y)\in I\}$, with basis $(\varphi f,
g\varphi^{-1})$. By the structure identity principle, $\varphi$ gives a
path $\mathrm{sip}(\varphi):A=B$.

\begin{lemma}[transport of structure]\label{lem:transport}
Transport along $\mathrm{ap}_{\RI}(\mathrm{sip}(\varphi))$ is
$I\mapsto\varphi[I]$.
\end{lemma}
\begin{proof}
By induction on isomorphisms it is enough to take $\varphi=1$, where both
sides are $I$: $\mathrm{sip}(1)$ is reflexivity, and $1[I]$ differs from $I$
only in proofs of propositions.
\end{proof}

\begin{lemma}[the triangle]\label{lem:triangle}
For $p\in\PGL$ and $q\in\Pn$, $\alpha(p)[\delta(q)]=\delta(\beta(p)q)$.
\end{lemma}
\begin{proof}
Both sides are right ideals, so it is enough to compare members, and we may
take $p=[G]$, $q=[x]$. With $H=G^{-1}$, $Y\in\alpha(p)[\delta[x]]$ means
$\Col_x(HYG)$. The identity $P(xr^T)Q=(Px)(Q^Tr)^T$ shows that $\Col_x(Z)$
implies $\Col_{Px}(PZQ)$, when $x$ has an invertible coordinate. Applied with $(G,H)$ and with $(H,G)$, this gives
$\Col_x(HYG)\iff\Col_{Gx}(Y)$.
\end{proof}

\begin{proposition}[sevb 3.3.1]\label{prop:331}
$\RI:\AZ\to\SB$ is an equivalence.
\end{proposition}
\begin{proof}
\emph{Sheaves.} $\SB$ is a $\Sigma$ over the type of sheaves (H3) of
sheaves. $\AZ$ is a retract of a $\Sigma$ over the type of sheaves $X$ of
the algebra structures on $X$ and of $\trT{-=\M}$. Algebra structures on a
sheaf form a sheaf: the operations are maps into $X$, and the laws, as
well as ``$X$ is a set'', are built from identity types of $X$.

\emph{Connected and pointed.} $\trT{A=\M}$ and $\trT{A'=\M}$ give
$\trT{A=A'}$, the other components being propositions; likewise for $\SB$.
The base points are $\M$ and $\Pn$, and $\RI(\M)=\Pn$ by
Lemma~\ref{lem:delta}.

\emph{Loops.} Since the other components are propositions, loops at $\M$
in $\AZ$ are loops in algebras, hence automorphisms of $\M$ by the
structure identity principle, hence $\PGL$ by $\alpha$. This gives an
equivalence $h:\PGL\simeq\Omega(\AZ,\M)$, $p\mapsto\mathrm{sip}(\alpha p)$.
Loops at $\Pn$ in $\SB$ are by univalence the automorphisms of $\Pn$, via
$k(\ell)=$ transport along $\ell$. For $p\in\PGL$, $\Omega\RI(h\,p)$ is the loop
$\mathrm{ap}_{\RI}(\mathrm{sip}(\alpha p))$ conjugated by the path
$\RI(\M)=\Pn$ of Lemma~\ref{lem:delta}, whose transport is $\delta^{-1}$. By
Lemmas~\ref{lem:transport} and~\ref{lem:triangle}, transport along it is
$q\mapsto\delta^{-1}(\alpha(p)[\delta(q)])=\beta(p)q$. So
$k\circ\Omega\RI\circ h=\beta$, and $\Omega\RI$ is an equivalence. Conclude
by Lemma~\ref{lem:deloop}.
\end{proof}

\section{Ch\^atelet's theorem}

\begin{lemma}[sevb 3.4.1]\label{lem:341}
Let $A:\AZ$, $I:\RI(A)$ and $(f,g)$ a basis of $I$. Then
$\sigma(a)_{ij}=g(f(e_i)\,a)_j$ is an algebra isomorphism $A\simeq\M$.
\end{lemma}
$\sigma(a)$ is the transpose of the matrix of $v\mapsto va$ on $I$ in the
basis $f$; so this is ``$A\simeq\mathrm{End}(I)^{op}\simeq\M$''.
\begin{proof}
That $\sigma$ is a homomorphism with contractible fibres is a sheaf
proposition (the carriers are sheaves), so by $\trT{A=\M}$ we may take
$A=\M$. Then $I=\delta[x]$ with $x_k$ invertible, by
Lemma~\ref{lem:delta}, and $f(e_i)=xy_i^T$. Put $D_{jk}=g(xe_k^T)_j$. Since
$w\mapsto g(xw^T)_j$ is linear, $\sigma(M)=YMZ$ with $Y_{im}=(y_i)_m$ and
$Z_{kj}=D_{jk}$. Moreover $YZ=1$, since $(YZ)_{ij}=g(xy_i^T)_j=g(f(e_i))_j$.
By Lemma~\ref{lem:K2}, $ZY=1$, so $\sigma$ is conjugation, an automorphism.
\end{proof}

\begin{proof}[Proof of Theorem~\ref{thm:main}]
The goal is a proposition. By Proposition~\ref{prop:331}, $X=\RI(A)$ for
some $A:\AZ$, so a point of $X$ is a right ideal $I:\RI(A)$. By (H4) we may
take a basis of $I$. Lemma~\ref{lem:341} gives $A\simeq\M$, hence $A=\M$,
and $X=\RI(A)=\RI(\M)=\Pn$ by Lemma~\ref{lem:delta}.
\end{proof}

\section{The \'etale topology}\label{sec:etale}

For $f=X^{d+1}+a_dX^d+\dots+a_0$ monic of positive degree, let $\Roots(f)$ be
the type of $x:R$ with $f(x)=0$. It is $\mathrm{Spec}(R[X]/(f))$, and
$R[X]/(f)$ is free with basis $1,X,\dots,X^d$. The \emph{\'etale topology} is
given by the propositions $G_f=\tr{\Roots(f)}$, indexed by the coefficients
$(d,a_0,\dots,a_d)$ (sevb 1.1.1--1.1.2).

\begin{lemma}\label{lem:sets}
A set $X$ is a sheaf if its identity types are sheaves and every map
$G_f\to X$ is merely constant. A set that is merely a sheaf is a sheaf.
\end{lemma}
\begin{proof}
For $h:G_f\to X$, ``$h$ is constant with value $x$'' determines $x$, since
identity types are sheaves; so the mere value is a value. The second claim
follows from the first.
\end{proof}

\begin{lemma}[sevb 1.1.3, 1.1.4]\label{lem:Rsheaf}
$R$ is a sheaf, and so is every affine scheme.
\end{lemma}
\begin{proof}
Let $B=R[X]/(f)$. By synthetic quasi-coherence, $B\to(\Roots(f)\to R)$ is a
bijection. Given $h:G_f\to R$, let $a\in B$ correspond to $x\mapsto h(x)$.
For a root $x$, this function is constant with value $h(x)$, so $a=h(x)\cdot1$.
Hence the constant coefficient $r$ of $a$ satisfies $r=h(x)$ for every
root, i.e.\ $r=h(g)$ for all $g:G_f$. Uniqueness: if $G_f\to r=s$, then
$r\cdot1$ and $s\cdot1$ agree at every root, so they are equal in $B$, and
$r=s$ since $1$ is part of a basis. Affine schemes are solution sets of
equations, so they are sheaves by Lemma~\ref{lem:closure}.
\end{proof}

\begin{lemma}\label{lem:opens}
Open propositions are sheaves.
\end{lemma}
\begin{proof}
$\Roots(f)$ is not empty: otherwise $\mathrm{Spec}(B)$ is empty, so $1=0$ in
$B$, while $R\to B$ is injective and $1\neq0$ in $R$. So $G_f\to U$ gives
$\neg\neg U$, and open propositions are double-negation stable.
\end{proof}

\begin{lemma}[sevb 1.2.2]\label{lem:compact}
For $U$ an open subset of $\Roots(f)$, the proposition $\forall
x.\,U(x)$ is open.
\end{lemma}
\begin{proof}
The goal is a proposition, and merely there are $h_1,\dots,h_m\in B$ such
that $U(x)$ holds exactly when some $h_s(x)$ is invertible. Then $\forall x.\,U(x)$ holds exactly when the
$h_s$ generate the unit ideal of $B$, i.e.\ when the $R$-linear map
$B^m\to B$, $(b_s)\mapsto\sum_sh_sb_s$, is surjective. In the basis of $B$
this is a matrix over $R$. Surjectivity of a matrix with $k+1$ columns is
equivalent to: for some row $j$ the last entry is invertible and the Schur
complement at that entry is surjective. By induction on $k$ this is a
finite disjunction of conjunctions of open propositions, hence open.
\end{proof}

This replaces sevb's argument through a closed subscheme of $\mathbb P^1$.

\begin{lemma}\label{lem:opentype}
The type of open propositions is a sheaf.
\end{lemma}
\begin{proof}
Equality of open propositions is logical equivalence, which is a sheaf by
Lemma~\ref{lem:opens}. Given $h:G_f\to\mathrm{Open}$, the proposition
$\forall y:\Roots(f).\,h(y)$ is open by Lemma~\ref{lem:compact}, and it is
equivalent to $h(g)$ for every $g:G_f$, since $G_f$ is a proposition.
Conclude by Lemma~\ref{lem:sets}.
\end{proof}

\begin{proposition}[sevb 1.2.1, 1.2.3]\label{prop:schemes}
Every scheme is a sheaf.
\end{proposition}
\begin{proof}
By Lemma~\ref{lem:sets} we may fix a finite cover of $X$ by open affine
charts $U_s$. \emph{Identity types.} Given $G_f\to x=y$, merely $x$ is in
some chart $U$; ``$y\in U$'' is open, hence a sheaf, so $y\in U$, and $U$ is
affine, hence a sheaf (Lemma~\ref{lem:Rsheaf}), so $x=y$. \emph{Maps.} Let
$h:G_f\to X$. For each $s$, $g\mapsto(h(g)\in U_s)$ is a map into open
propositions, hence constant with some value $V_s$
(Lemma~\ref{lem:opentype}). Since the charts cover, $G_f\to\bigvee_sV_s$, a
finite disjunction of opens, hence a sheaf, so $\bigvee_sV_s$ holds. For
$s$ with $V_s$, $h$ lands in the sheaf $U_s$, so it is constant.
\end{proof}

\begin{corollary}
For the \'etale topology, Theorem~\ref{thm:main} holds under (H2) and (H3).
\end{corollary}

(H2) and (H3) hold for localisation at a family of propositions
\cite{RSS20}; they are not formalised here. (H4) is the use made of sevb's
Proposition~1.3.5; it is proved in \cite{descent} for algebras with
$\trT{A=\M}$, which is all that the proof of Theorem~\ref{thm:main}
uses, without any theory of flatness.

\section{Correspondence with the formalisation}

Everything is in \texttt{sag/chatelet}: about 4000 lines in 46
files, on top of the library for synthetic algebraic geometry and
projective space, which is included in the same directory; the descent files
\texttt{Descent*} are described in \cite{descent}. The tree of 393 files
checks in 11~s in one process. Every file checks in under 0.2~s, except
\texttt{Prop331Str} (2~s). The hypotheses are module parameters, never
postulates:
\texttt{sag} (the SAG axioms), \texttt{sh} (H2 and H3: the record
\texttt{SheafClosure.Ops}, with field \texttt{lex}), \texttt{R-modal} (H1),
\texttt{descends} (H4), the last one discharged for the \'etale topology
by \texttt{ChateletEtaleDescent}. Functional extensionality is derived from
univalence. Large types are handled by universe levels: $\RI(A)$ and $\AZ$
are one level up, $\SB$ two levels up.

\medskip
{\small
\begin{longtable}{p{4.6cm}p{8.4cm}}
\textbf{Statement} & \textbf{File and name}\\ \hline
Sheaves, $\trT{-}$ (H2), (H3) & \texttt{SheafTopology.isModal},
  \texttt{SheafClosure.Ops}\\
Lemma \ref{lem:closure} & \texttt{SheafLemmas}: \texttt{modal-Pi},
  \texttt{modal-Sigma}, \texttt{modal-Id}, \texttt{modal-retract},
  \texttt{modal-of-sec}\\
Lemma \ref{lem:propT} & \texttt{AzumayaProbe}: \texttt{PropT},
  \texttt{modal-PropT}, \texttt{propT-ext}\\
Lemma \ref{lem:deloop} & \texttt{Deloop.Lemma.theorem},
  \texttt{DeloopLevels.deloop}\\
Lemmas \ref{lem:K1}, \ref{lem:K2} & \texttt{MatrixLeftInverse}:
  \texttt{pivot}, \texttt{k2}\\
$\M$, $E_{ij}$ & \texttt{MnAlgebra.algebra}, \texttt{MnElementary.E}\\
Proposition \ref{prop:sn} & \texttt{SkolemNoether.Endo.theorem}\\
Corollary \ref{cor:alpha} & \texttt{SkolemNoetherPGL}: \texttt{alpha},
  \texttt{equivalence}, \texttt{equivalence-iso}\\
$\beta$ & \texttt{AutPGL.aut-of}, \texttt{corollary-5-6}; instances in
  \texttt{AutPn}\\
Definition \ref{def:ri}, $\AZ$, $\SB$ & \texttt{AzumayaProbe}:
  \texttt{Ideal}, \texttt{Basis}, \texttt{Retraction}, \texttt{RI},
  \texttt{RI-modal}, \texttt{Types.AZ}, \texttt{Types.SB}\\
Lemma \ref{lem:delta}, sevb 2.4.3 & \texttt{Delta}: \texttt{Col},
  \texttt{Memb}, \texttt{delta-map}, \texttt{inj}, \texttt{fib-point},
  \texttt{delta}, \texttt{lemma-2-4-3}\\
Structure identity principle & \texttt{AlgSIP.sip}, \texttt{AlgSIPInd.iso-ind}\\
Lemma \ref{lem:transport} & \texttt{TriangleRI}: \texttt{RIm.map},
  \texttt{transport}\\
Lemma \ref{lem:triangle} & \texttt{Triangle}: \texttt{OfA.tri-RI},
  \texttt{tri}, \texttt{omega}\\
Proposition \ref{prop:331} & \texttt{Prop331.prop-3-3-1}; algebra
  structures: \texttt{Prop331Str.str-modal}; isEquiv is a proposition:
  \texttt{Prop331Paths.isEquiv-prop}\\
Lemma \ref{lem:341} & \texttt{Lemma341.lemma-3-4-1}\\
Theorem \ref{thm:main} & \texttt{Theorem342.theorem-3-4-2}\\
\'Etale topology & \texttt{EtaleTopology}: \texttt{I}, \texttt{Root},
  \texttt{G}\\
Lemma \ref{lem:sets} & \texttt{EtaleModal}: \texttt{Separated.separated-factor},
  \texttt{Mere.modal}\\
Lemma \ref{lem:Rsheaf} & \texttt{EtaleRing}: \texttt{R-local},
  \texttt{affine-local}; \texttt{EtaleRoots.spec-root}\\
Lemma \ref{lem:opens} & \texttt{EtaleOpen.open-local};
  \texttt{EtaleRoots}: \texttt{not-not-root}\\
Lemma \ref{lem:compact} & \texttt{EtaleCompact.At.compact};
  \texttt{EtaleIdealOpen.unimodular-open}, \texttt{EtaleSurj.surj-open}\\
Lemma \ref{lem:opentype} & \texttt{EtaleOpenType.opens-local}\\
Proposition \ref{prop:schemes} & \texttt{EtaleScheme.etale-scheme}\\
Corollary (\'etale) & \texttt{ChateletEtale.At.chatelet}\\
Corollary, (H4) proved & \texttt{ChateletEtaleDescent.At.chatelet}\\
\end{longtable}
}

\section{Remarks on sevb}
\sloppy

The numbering is that of sevb. Sections 2.2, 2.3 and 2.5 of sevb have
not been formalised; Section 1.3 is replaced by \cite{descent}.

\subsection*{Gaps}
\begin{enumerate}
\item \emph{Lemma 2.4.2 (surjectivity) and Proposition 3.3.1: ``same
  dimension, hence equal''.} Over a ring a free submodule of the same rank
  need not be the whole module: $(t)\subseteq R$ is free of rank one when
  $t$ is a non-zero-divisor. The step holds when the inclusion has a left
  inverse, by Lemma~\ref{lem:K2}. So $\mathrm{Gr}_k$ should consist of
  split subspaces (frames with a left inverse), which is what ``matrices of
  rank $k$'' in the proof of Lemma 2.1.4 should mean. In Proposition 3.3.1
  no dimension argument is needed: the inclusion applied to $P^{-1}$ gives
  the other inclusion.
\item \emph{Theorem 3.4.2, ``we merely have that $I=R^{n+1}$''.} From
  $I:\RI(A)$ one only has $\trT{I=R^{n+1}}$. Passing to $\tr{A=M_{n+1}(R)}$
  uses that ``finite free'' is a sheaf proposition (Proposition 1.3.5, or the
  hypothesis that the type of finite free modules is a sheaf). This is the
  only use of descent in the main line; it should be explicit.
\item \emph{Definition 1.1.1, Remark 1.1.2: the degree of $g$ must be
  positive.} For $g=1$, $\mathrm{Spec}(R[X]/g)$ is empty, and only
  contractible types are local. The proof of Lemma 1.1.3 fails for
  $n=0$ (the equaliser of $R^0\rightrightarrows R^0\otimes R^0$ is $0$), and
  ``$\neg\neg\mathrm{Spec}(R[X]/g)$'' in Proposition 1.2.3 is false. For
  positive degree it holds since an empty spectrum gives $1=0$ in
  $R[X]/g$, while $R\to R[X]/g$ is injective; this deserves a line.
\end{enumerate}

\subsection*{Imprecisions and simplifications}
\begin{itemize}
\item \emph{Lemma 2.1.5, ``we can proceed as if $R$ was a field''.}
  Extending a basis of $I$ to a basis of $A$ presupposes a complement. With
  a left inverse $g$ of the frame $f$, ``$x\in I$'' is the closed condition
  $f(gx)=x$.
\item \emph{Lemma 3.2.2, ``from $V_0^{n+1}=R^{n+1}$ we get $\tr{V_0=R}$''.}
  This uses uniqueness of rank, which is not stated, after Lemma 3.2.1. The
  detour through projective modules can be avoided
  (Proposition~\ref{prop:sn}): by locality one finds $j$ and $w$ in the
  image of $e_{00}$ with a linear form $\mu$, $\mu(w)=1$; then
  $c\mapsto\sum_kc_ke_{k0}w$ has an explicit left inverse, so it is
  invertible by Lemma~\ref{lem:K2}, and it intertwines $e_{ij}$ and
  $E_{ij}$. Neither Lemma 3.2.1 nor ranks are needed.
\item \emph{Proposition 3.2.3 and Lemma 2.4.2.} ``$\lambda\neq0$'' should be
  ``$\lambda$ a unit'' (equivalent in synthetic algebraic geometry, but
  this should be said).
\item \emph{Proposition 3.3.1.} The formula for $\delta(Y)$ is garbled ($Y$
  versus $X$, $M_n$, and the products do not type check). With
  $\delta[x]=\{xy^T\}$ the triangle is the equality of sets
  $P\,\delta[x]\,P^{-1}=\delta[Px]$, since $P(xy^T)P^{-1}=(Px)(P^{-T}y)^T$.
  ``$\mathrm{ap}_{\RI}(\alpha(P))$'' should be transport along
  $\mathrm{ap}_{\RI}(\mathrm{sip}(\alpha P))$; that this is $I\mapsto
  \varphi[I]$ is a lemma (Lemma~\ref{lem:transport}). Lemma 3.1.2 needs
  $\AZ$ and $\SB$ to be connected sheaves, which is not checked: $\SB$ is a
  sheaf by lexness, and $\AZ$ because algebra structures on a sheaf form a
  sheaf. The identifications $\Omega\AZ=\Aut(M_{n+1}(R))$ (structure identity
  principle) and $\Omega\SB=\Aut(\Pn)$ (univalence) are implicit.
\item \emph{Lemma 3.4.1.} ``Sends $M$ to $\delta^{-1}\circ M\circ\delta$''
  should be $\theta^{-1}\circ M^T\circ\theta$.
\item \emph{Lemma 1.1.3.} A shorter proof: by synthetic quasi-coherence
  $R[X]/g\simeq(\mathrm{Spec}(R[X]/g)\to R)$, and a map $\tr{\mathrm{Spec}}\to
  R$ corresponds to an element equal to $h(x)\cdot1$, whose constant
  coefficient is the value (Lemma~\ref{lem:Rsheaf}).
\item \emph{Lemma 1.2.2.} An algebraic proof, without $\mathbb P^1$, is
  Lemma~\ref{lem:compact}.
\item \emph{Hypotheses of Section 2.} With right ideals valued in sheaf
  propositions, the main line uses only: $R$ is a sheaf, lexness, and
  descent of finite freeness. That $\Pn$ is a sheaf then follows from
  Lemma~\ref{lem:delta}. ``Schemes are sheaves'' is used only for Remark
  3.3.2.
\item \emph{Lemma 2.5.2.} ``We can assume $x_0\neq0$'' needs a case split
  on which coordinate of the point is invertible, not only
  $C(a,b)=C(1/a,-b/a)$.
\end{itemize}

\subsection*{Typos}
$M_n(R)$ for $M_{n+1}(R)$ in Lemma 2.4.2 (twice), Proposition 3.3.1 (the
triangle and $\delta(Y)$) and Lemma 3.4.1; Lemma 1.2.1, ``$f^{-1}(U_i)$ is
an in $P$'' (open); Lemma 2.5.2, ``$ay_0+bz_0=1$'' for $ay_0^2+bz_0^2=1$, and
``$\phi$ inverse to $\phi$'' for ``inverse to $\psi$''; ``generalise of'',
``over over'', ``comuutative'' (introduction), ``when when'' (2.4.2), ``an
well-known'' (3.2), ``Let us proof'' (2.5.2), ``conclude than'' (3.3.2).

\begin{thebibliography}{9}
\bibitem{CCH24} F.~Cherubini, T.~Coquand, M.~Hutzler, \emph{A foundation for
synthetic algebraic geometry}, 2024.
\bibitem{Che+24} F.~Cherubini, T.~Coquand, M.~Hutzler, D.~W\"arn,
\emph{Projective space in synthetic algebraic geometry}, 2024.
\bibitem{descent} \emph{Descent of finite freeness}, note,
\texttt{doc/chatelet-descent.pdf}, 2026.
\bibitem{RSS20} E.~Rijke, M.~Shulman, B.~Spitters, \emph{Modalities in
homotopy type theory}, LMCS 2020.
\bibitem{sevb} \emph{Ch\^atelet's theorem in synthetic algebraic geometry},
note, October 2026.
\end{thebibliography}

\end{document}
